\documentclass[11pt,twoside,openright]{book}

\usepackage[T1]{fontenc}
\usepackage[utf8]{inputenc}
\usepackage{amsmath,amssymb,amsthm,mathtools,bm}
\usepackage{mathpazo}
\usepackage{microtype}
\usepackage[inner=1.35in,outer=1.1in,top=1.25in,bottom=1.35in]{geometry}
\usepackage{booktabs}
\usepackage{fancyhdr}
\usepackage[hidelinks]{hyperref}
\usepackage{cleveref}

\setlength{\parindent}{1.4em}
\setlength{\parskip}{0pt}
\linespread{1.07}
\emergencystretch=3em

\pagestyle{fancy}
\fancyhf{}
\fancyhead[LE]{\small\thepage\quad\textsc{A Cinema of Individual Editions}}
\fancyhead[RO]{\small\nouppercase{\leftmark}\quad\thepage}
\renewcommand{\headrulewidth}{0.3pt}
\fancypagestyle{plain}{\fancyhf{}\fancyfoot[C]{\small\thepage}\renewcommand{\headrulewidth}{0pt}}

\theoremstyle{plain}
\newtheorem{theorem}{Theorem}[chapter]
\newtheorem{proposition}[theorem]{Proposition}
\newtheorem{lemma}[theorem]{Lemma}
\newtheorem{corollary}[theorem]{Corollary}
\theoremstyle{definition}
\newtheorem{definition}[theorem]{Definition}
\newtheorem{example}[theorem]{Example}
\theoremstyle{remark}
\newtheorem{remark}[theorem]{Remark}

\newcommand{\Capac}{\mathrm{Cap}}
\newcommand{\Rob}{\mathrm{Rob}}
\newcommand{\Inv}{\mathrm{Inv}}
\DeclareMathOperator{\loss}{loss}

\title{\Huge\bfseries A Cinema of Individual Editions\\[1.2ex]
\large Chapter Eight: Amendments, Footprints, and Correction Events\\[0.4ex]
\large Working Draft}
\author{Flyxion\\[0.3ex]\small Independent Researcher}
\date{November 2026}

\begin{document}

\frontmatter
\maketitle
\tableofcontents
\mainmatter
\setcounter{chapter}{7}

\chapter{Amendments, Footprints, and Correction Events}
\label{ch:amend}

Chapter~7 ends with an upheld report. Upheld reports have two possible outcomes. A report may show that an edition violated a rule already in force, so that the assembly or the assigned take was at fault and the rules do not change. Or it may supply evidence for an amendment, so that the rules were incomplete and must be extended. Violations lead to repair, and this chapter is about the other outcome: how a rule is derived from one instance, how far it reaches, how the correction is made as a single event, when it cannot be made, and what a production can do when it cannot.

The chapter joins the earlier ones. Chapter~5 gave the footprint of a defect, Chapter~9 the capacity cost of an amendment, and Chapter~10 the repair of the viewers it invalidates. The correction event of this chapter is the unit that carries all three at one moment, so that the ledger never shows rules that the assignments do not meet. Production C is again the stress case, and it supplies the chapter's last result: where replacement footage is shot matters for repair even when it makes no difference to capacity.

\section{From One Instance to a Rule}\label{sec:general}

An upheld report names a strip, a slot, and a stale class, and it is one instance of a defect. The rule it supports is a generalization of the instance, and there are three natural ones, which are the levels of Chapter~5: the interval of the defect, every showing of the entity, and every entity of the type. They cover different sets of stale classes. \Cref{tab:cover} gives, for each first report in each production, how many of the pool's stale classes each level removes and the capacity that remains.

\begin{table}[ht]
\centering
\begin{tabular}{@{}llccccccc@{}}
\toprule
Prod. & First report & \multicolumn{3}{c}{Stale classes removed} & \multicolumn{3}{c}{Capacity remaining}\\
\cmidrule(lr){3-5}\cmidrule(lr){6-8}
 & (strip, slot) & Interval & Entity & Type & Interval & Entity & Type\\
\midrule
A (4 stale) & $a$, 2 & 2 & 2 & 4 & 351 & 351 & 238\\
 & $c$, 4 & 1 & 1 & 4 & 422 & 422 & 238\\
 & $b$, 5 & 1 & 1 & 4 & 404 & 404 & 238\\
B (7 stale) & $a$, 2 & 2 & 2 & 7 & 200 & 200 & 96\\
 & $b$, 3 & 2 & 2 & 7 & 288 & 288 & 96\\
 & $c$, 5 & 3 & 3 & 7 & 176 & 176 & 96\\
C (2 stale) & $a$, 2 & 1 & 1 & 2 & 1 & 1 & 1\\
 & $c$, 5 & 1 & 1 & 2 & 2 & 2 & 1\\
\bottomrule
\end{tabular}
\caption{For each production and each first report, the number of the pool's stale classes removed by the interval, entity, and type levels, and the capacity that remains. The interval and entity levels coincide in every row because no stale class shows the reported strip in the other interval.}\label{tab:cover}
\end{table}

The table shows what a single instance can and cannot justify. The interval level removes exactly the stale classes that the instance represents, those of the same strip in the same persistence interval, and Proposition~5.4 of Chapter~5 proves it sound and complete for that interval. The type level removes every stale class of the pool, which is more than the instance supports in every row: after the first report on strip $a$ in production A, for instance, it removes two stale classes that the report does not represent. The type-level rule is justified only if the production has grounds to think the defect belongs to the whole type, and one instance does not supply them.

\begin{definition}[Supported rule]\label{def:supported}
Given a set $\mathcal R$ of upheld reports, the \emph{supported rule} is the union of the interval-level amendments of its reports. It can be raised to the entity level or the type level only on further evidence, namely reports of the same defect in other intervals of the entity or on other entities of the type.
\end{definition}

The audience of production A reveals every pair of strip and interval that its pool contains, so the union of the interval amendments of its reports has the capacity of $G$ on that pool, $238$. That is a fact about this pool and audience, and the production could not have known it beforehand. Production C's audience reveals one of the two pairs in its pool, and the supported rule from its reports is the interval amendment on strip $a$ alone, which leaves the defect on strip $c$ in the pool for a later report.

\section{The Order in Which Corrections Arrive}\label{sec:order}

Reports arrive in an order that no one controls, and each upheld pair of strip and interval can trigger its own amendment. The order does not change where the production ends and it changes the path.

\begin{proposition}[Order independence]\label{prop:order}
Let amendments $X_1',\dots,X_n'$ be applied in any order. The admissible set after all of them is $\bigcap_iX'_i$ whatever the order, and so is its capacity. The capacity along the way decreases from step to step, so the minimum capacity along any order is the final capacity, and the promise to an audience of $U$ is kept at every step if and only if $\Cap(\bigcap_iX_i')\ge U$.
\end{proposition}

\begin{proof}
Removing realizations is commutative, since each amendment is a subset and applying one after another gives the intersection. Each removal can only reduce capacity, so the sequence of capacities is nonincreasing and ends at the capacity of the intersection.
\end{proof}

The intermediate values differ. For production A the three interval amendments, on strip $a$ after its event, strip $b$ after its event, and strip $c$ before its event, can arrive in six orders. Every order ends at $238$. The first step leaves $351$, $404$, or $422$ depending on which strip is reported first, and the middle step leaves between $286$ and $332$. For production B the first step leaves $200$, $288$, or $176$, the middle step between $128$ and $168$, and the end is $96$. The paths are in \Cref{tab:paths}. The differences matter to what follows, because the repair radii of Chapter~10 are computed against the pool at each step and not against the final pool.

\begin{table}[ht]
\centering
\begin{tabular}{@{}llccc@{}}
\toprule
Prod. & Order of reports & After first & After second & After third\\
\midrule
A & $a, b, c$ & 351 & 286 & 238\\
A & $a, c, b$ & 351 & 303 & 238\\
A & $b, a, c$ & 404 & 286 & 238\\
A & $b, c, a$ & 404 & 332 & 238\\
A & $c, a, b$ & 422 & 303 & 238\\
A & $c, b, a$ & 422 & 332 & 238\\
B & $a, b, c$ & 200 & 168 & 96\\
B & $a, c, b$ & 200 & 128 & 96\\
B & $b, a, c$ & 288 & 168 & 96\\
B & $b, c, a$ & 288 & 144 & 96\\
B & $c, a, b$ & 176 & 128 & 96\\
B & $c, b, a$ & 176 & 144 & 96\\
\bottomrule
\end{tabular}
\caption{Capacity after each of three interval-level amendments, in every order. Premiere capacities are $494$ for A and $320$ for B.}\label{tab:paths}
\end{table}

\section{The Correction Event}\label{sec:event}

A correction is formed as one event, as the first chapter requires, with three parts at one effective time: the amendment of the rules, any footage added to the pool, and the repairs of the allocation. The reason for atomicity is that the ledger is the evidence of adjudication, and a ledger that shows the new rules beside the old assignments records editions that its own rules call inadmissible.

\begin{definition}[Correction event]\label{def:event}
A \emph{correction event} at effective time $\tau$ is a triple of an amendment $X'\subseteq X$, an addition $\Delta T$ of takes to the pool, and a repair $a\mapsto a'$ of the allocation. It is \emph{formable} if the pool after the amendment and the addition admits an allocation for the audience, namely $\Cap(X'\cup\Delta X)\ge U$, and the repair keeps every moved viewer within the permitted budget.
\end{definition}

The cost of not forming events atomically can be counted. Suppose the three amendments of production A are adopted in turn, each repaired afterward as in Chapter~10. If each rule enters the ledger before its repairs, the viewers whose editions the rule has just made inadmissible number $16$, then $15$, then $6$, a total of $37$ viewer-revisions on inadmissible editions. For production B the numbers are $8$, $3$, and $15$, a total of $26$. With atomic events the count is zero and the admissibility invariant of Chapter~1, that every assignment in the ledger meets the rules of its revision, holds after every event, which was checked for both productions.

\subsection{When the Event Cannot Be Formed}

An event can fail in two ways. The capacity condition can fail, as for production C under the agreement rule: the capacity is $1$ for an audience of $2$, and no repair of any radius exists. The budget condition can fail with capacity to spare. For A and B under the full agreement rule, the survivors-fixed repair needs radius $2$, and with a budget of $1$ no repair exists although $238$ and $96$ class tuples remain. The event of production A is formable at a budget of $2$ or $3$ and not at $1$, and the same holds for B.

When an event cannot be formed there are four things the production can do: hold the amendment and keep the defect in the pool, which breaks the promise of Chapter~1 to the viewers who would see it; raise the budget, which disturbs viewers more; add footage; or relax a rule that the history permits. The last two are the subject of the next two sections.

\section{Replacement Footage}\label{sec:footage}

Footage that is added to the pool restores capacity, and where it is added decides who can use it. For production C under $G$, cloning its admissible class at any one slot raises the capacity from $1$ to $2$, which is exactly the audience. The two class tuples that result differ at the cloned slot. The viewer who held the stale edition must move to the new tuple, and the distance depends on where the clone is. At slot $2$, the slot of the stale class, the viewer moves one slot, from the stale class to its correct replacement, and the repair radius is $1$. At every other slot the radius is $2$: the viewer moves from the stale class at slot $2$ and from the original class to the clone at the other slot. \Cref{tab:clone} gives all five.

\begin{table}[ht]
\centering
\begin{tabular}{@{}lccccc@{}}
\toprule
Clone at slot & 1 & 2 & 3 & 4 & 5\\
\midrule
Capacity & 2 & 2 & 2 & 2 & 2\\
Repair radius & 2 & 1 & 2 & 2 & 2\\
\bottomrule
\end{tabular}
\caption{Production C after $G$: the effect of adding one clone of the admissible class at each slot, for the two viewers of the premiere.}\label{tab:clone}
\end{table}

Capacity alone cannot choose between the five placements, and the repair radius can. The footage that fixes a defect in place is the footage that costs the viewer least, which is a reason to reshoot the slot that held the error before shooting anywhere else. This is a statement about a single instance of C, and it need not hold in general, since a clone at a slot of heavy traffic can add many more tuples, as Chapter~9 showed for production A.

\section{Relaxing a Rule That the History Permits}\label{sec:relax}

A rule can be relaxed when the story world admits the depiction it forbids. The direction rule forbids a cut between shots with opposite screen directions, and a story in which the character turns, or crosses the line of action on camera, permits the cut. Relaxing the rule at one edge enlarges the admissible set. \Cref{tab:relax} gives the capacity after $G$ when the direction rule is dropped at each of the four edges of the scene.

\begin{table}[ht]
\centering
\begin{tabular}{@{}lrrrrr@{}}
\toprule
Production & After $G$ & Edge (1,2) & Edge (2,3) & Edge (3,4) & Edge (4,5)\\
\midrule
A & 238 & 310 & 364 & 270 & 260\\
B & 96 & 200 & 200 & 192 & 192\\
C & 1 & 2 & 1 & 1 & 1\\
\bottomrule
\end{tabular}
\caption{Capacity after the agreement rule when the direction rule is relaxed at one edge.}\label{tab:relax}
\end{table}

For production C, relaxing the first edge alone restores a second edition: slot $1$ may be either of its two classes, since the rule no longer links it to slot $2$. The repair is not cheap. The viewer who held the stale edition moves from the stale class at slot $2$ to the correct class and from the first slot's class to the other, a radius of $2$.

The model cannot say which edges the history permits. That belongs to the verdict of Chapter~7: the third step compares the verified appearance with the script history, and the same history decides whether a relaxation is licensed. A relaxation that the history does not license is not a correction but an error of a different kind, since it admits editions that depict a different story.

\section{One Process}\label{sec:process}

Footage, relaxation, the level of the rule, the budget, and the order of reports are parts of one decision. For production C after the first defect report, doing nothing keeps an edition that the viewer's report has shown to be wrong. The agreement rule forbids it and leaves the promise unkept. A clone at slot $2$ keeps the promise at radius $1$. A clone elsewhere keeps it at radius $2$. A relaxation at the first edge keeps it at radius $2$, at no storage cost and at the cost of a change to the film's grammar. None of these is forced by the mathematics. The mathematics says which are feasible and what each costs, in tuples, in disturbance, and in footage.

The scripts \texttt{lawnmodel.py} and \texttt{verify\_ch08.py} reproduce each figure of this chapter.

\begin{thebibliography}{99}

\bibitem{flyxion2026gf}
Flyxion.
\newblock \emph{Generative Film: Categories, Sheaves, and Types for Cinema}.
\newblock Unpublished monograph, October 2026.

\end{thebibliography}

\end{document}
